% UC3M / Instituto de Salud Carlos III presentation template -- short
% English deck. The theme contains no text strings of its own, so the only
% language-dependent line is the babel one below; everything else is
% identical to the Spanish demo. This file only shows one frame of each
% kind (two columns, normal, two rows, table); see main.tex for the full
% Spanish showcase of every layout and feature.
%
% The content is a made-up example (anomaly detection in sensor time
% series); the figures do not come from any real experiment.
%
% Build:  make      (latexmk -lualatex, two passes for the outline frames;
%                    produces build/main-en.pdf along with main.pdf)
\documentclass[aspectratio=169,9pt]{beamer}
\usetheme{uc3m}

\usepackage[english]{babel}
\usepackage{booktabs}

\title{Anomaly detection in sensor time series}
\date{September 2026}
\author{Raúl Aguilar Arroyo\\Name Surname\\Name Surname}

\begin{document}

\begin{frame}[plain,noframenumbering]
  \titlepage
\end{frame}

% ==========================================================================
%  TWO-COLUMN SLIDE
% ==========================================================================
\section{Two columns}

\begin{frame}{Heading and several paragraphs}
\begin{twocols}
  \panelheading{Background}

  Classical anomaly-detection methods for time series fit a model of the
  signal and flag as anomalous the readings that stray from the
  prediction.

  Their main cost is fitting the model, which has to be repeated for every
  sensor and revisited whenever operating conditions change, with a
  typical calibration window of \textbf{[t-30, t-1]}.
\nextcol
  \panelheading{This work}

  We check whether using \textbf{only} a moving average and a threshold on
  the residual recovers that same detection capability.

  The detector is a sliding-window \emph{z-score} with the same time
  windows and a threshold set by cross-validation over partitions of the
  historical data.
\end{twocols}
\end{frame}

% ==========================================================================
%  NORMAL SLIDE
% ==========================================================================
\section{Normal}

\begin{frame}{Text with a figure below}
  \panelheading{Daily signal}

  Daily reading of a synthetic sensor with its 15-day moving average,
  aggregated over the 12 sensors of the installation.

  \ucmfigure[0.65]{figs/demo-figura}
\end{frame}

\begin{frame}{Table}
  \panelheading{Estimates}

  \begin{tabular}{@{}lrrr@{}}
    \toprule
    Threshold     & Precision & Recall & $F_1$ \\
    \midrule
    $z > 2.0$     & 0.71  & 0.88  & 0.79 \\
    $z > 2.5$     & 0.82  & 0.79  & 0.80 \\
    $z > 3.0$     & 0.91  & 0.61  & 0.73 \\
    $z > 3.5$     & 0.95  & 0.42  & 0.58 \\
    \bottomrule
  \end{tabular}

  Estimated by cross-validation over the 312 partitions with labelled
  anomalies.
\end{frame}

% ==========================================================================
%  TWO-ROW SLIDE
% ==========================================================================
\section{Two rows}

\begin{frame}{One paragraph per row}
\begin{tworows}
  The sliding-window detector compares each reading with those in its
  immediate neighbourhood: the reference is computed on the series itself,
  so anything that varies slowly is absorbed by the moving average.
\nextrow
  This removes the need for an explicit model of the signal and cancels
  slow sensor drift, at the cost of missing anomalies that extend beyond
  the window size.
\end{tworows}
\end{frame}

\end{document}
